\documentclass[11pt]{article}
\usepackage[a4paper,margin=1in]{geometry}
\usepackage{fontspec}
\usepackage[boldfont=false]{xeCJK}
\setCJKmainfont{FandolSong-Regular.otf}
\usepackage{amsmath}
\usepackage{amssymb}
\usepackage{array}
\usepackage{booktabs}
\usepackage{graphicx}
\usepackage{adjustbox}
\usepackage{enumitem}
\setlist[itemize]{topsep=2pt,partopsep=0pt,parsep=0pt,itemsep=2pt,leftmargin=1.4em}
\setlist[enumerate]{topsep=2pt,partopsep=0pt,parsep=0pt,itemsep=2pt,leftmargin=1.6em}
\usepackage{parskip}
\usepackage{fvextra}
\fvset{breaklines=true,breakanywhere=true,fontsize=\small}
\setlength{\parindent}{0pt}

\begin{document}

\begin{center}
  {\LARGE\bfseries The price system and the microeconomy}\\[0.35em]
  {\large A Level Economics}
\end{center}
\vspace{0.6em}

\section*{Demand}

\textbf{demand} 需求 is the quantity of a good that buyers are willing and able to buy at each \textbf{price} 价格 in a period of time. It is not just a wish — buyers must also be able to pay.

The \textbf{law of demand} says: when the price rises, the \textbf{quantity demanded} 需求量 falls; when the price falls, the quantity demanded rises. Price and quantity move in opposite directions.

So the \textbf{demand curve} 需求曲线 slopes downward, from top-left to bottom-right.

\subsection*{What shifts the demand curve}

A change in the good's \textbf{own price} is shown as a \emph{movement along} the demand curve. A change in any other cause shifts the whole curve to a new position — this is a \textbf{shift} 移动 of demand. The main causes (the \textbf{determinants} 决定因素) are:

\begin{itemize}
\item \textbf{income} 收入 — for most goods, higher income raises demand.

\item the price of a \textbf{substitute} 替代品 — a good you can use instead (tea for coffee). If the substitute's price rises, demand for this good rises.

\item the price of a \textbf{complement} 互补品 — a good used together with this one (cars and petrol). If the complement's price rises, demand for this good falls.

\item tastes and fashion, the size of the population, and advertising.

\end{itemize}

\subsection*{Individual and market demand}

One \textbf{consumer} 消费者 has an individual demand. The \textbf{market demand} 市场需求 is found by adding up the quantity that all consumers want at each price (adding the curves sideways).

\subsection*{Marginal utility and demand}

\textbf{utility} 效用 is the satisfaction you get from consuming a good. \textbf{marginal utility} 边际效用 is the extra utility from one more unit.

The law of \textbf{diminishing marginal utility} 边际效用递减 says each extra unit gives less extra satisfaction than the one before. Because later units are worth less to you, you will only buy them at a lower price. This is why the demand curve slopes downward.

\section*{Supply}

\textbf{supply} 供给 is the quantity of a good that \textbf{producers} 生产者 are willing and able to sell at each price in a period of time.

The \textbf{law of supply} says: when the price rises, the \textbf{quantity supplied} 供给量 rises. Higher prices mean more \textbf{profit} 利润, so firms make more. So the \textbf{supply curve} 供给曲线 slopes upward.

\subsection*{What shifts the supply curve}

A change in the good's own price is a movement along the supply curve. Other causes shift the whole curve:

\begin{itemize}
\item \textbf{costs of production} 生产成本 — if wages, raw materials or rent get dearer, supply falls.

\item technology — better machines raise supply.

\item a \textbf{subsidy} 补贴 (government payment to producers) raises supply; an indirect tax lowers it.

\item the number of \textbf{firms} 企业 in the market, and weather (for farm goods).

\end{itemize}

\subsection*{Individual and market supply}

The \textbf{market supply} 市场供给 is found by adding up the quantity all firms will sell at each price.

\section*{Price elasticity of demand}

\textbf{price elasticity of demand} 需求价格弹性 (PED) measures how much the quantity demanded responds when the price changes.

\[
\text{PED} = \frac{\%\ \text{change in quantity demanded}}{\%\ \text{change in price}}
\]

PED is normally negative (because of the law of demand), but we usually look at the size and drop the minus sign.

\begin{center}
\adjustbox{max width=\linewidth}{%
\begin{tabular}{lll}
\toprule
\textbf{Value} & \textbf{Name} & \textbf{Meaning} \\
\midrule
more than 1 & \textbf{elastic} 富有弹性 & quantity responds a lot \\
less than 1 & \textbf{inelastic} 缺乏弹性 & quantity responds a little \\
equal to 1 & \textbf{unit elastic} 单位弹性 & quantity changes by the same \% as price \\
0 & \textbf{perfectly inelastic} 完全无弹性 & quantity does not change at all \\
infinity & \textbf{perfectly elastic} 完全有弹性 & any price rise sends demand to zero \\
\bottomrule
\end{tabular}}
\end{center}

\subsection*{What makes demand elastic or inelastic}

\begin{itemize}
\item the number and closeness of substitutes (more substitutes → more elastic).

\item whether the good is a \textbf{necessity} 必需品 (inelastic) or a \textbf{luxury} 奢侈品 (elastic).

\item the share of income spent on the good (a large share → more elastic).

\item time — demand is more elastic in the long run, when buyers can find other options.

\item habit or addiction makes demand inelastic.

\end{itemize}

\subsection*{PED and total revenue}

\textbf{total revenue} 总收益 is the money a firm earns: $\text{TR} = \text{price} \times \text{quantity}$. PED tells a firm what happens to revenue when it changes the price:

\begin{center}
\adjustbox{max width=\linewidth}{%
\begin{tabular}{lll}
\toprule
\textbf{If demand is…} & \textbf{price rises} & \textbf{price falls} \\
\midrule
inelastic & total revenue rises & total revenue falls \\
elastic & total revenue falls & total revenue rises \\
unit elastic & total revenue stays the same & total revenue stays the same \\
\bottomrule
\end{tabular}}
\end{center}

This is useful to firms (how to price for more revenue) and to \textbf{governments} 政府 (a tax on an inelastic good, like petrol, raises a lot of revenue without cutting sales much).

\section*{Income elasticity of demand}

\textbf{income elasticity of demand} 需求收入弹性 (YED) measures how much demand responds when income changes.

\[
\text{YED} = \frac{\%\ \text{change in quantity demanded}}{\%\ \text{change in income}}
\]

\begin{itemize}
\item a \textbf{normal good} 正常品 has positive YED: demand rises as income rises. Necessities have YED between 0 and 1; luxuries have YED above 1.

\item an \textbf{inferior good} 低档品 has negative YED: demand falls as income rises (people switch to better goods). An example is cheap instant noodles.

\end{itemize}

This helps firms plan: as incomes grow, demand for luxuries grows fastest.

\section*{Cross elasticity of demand}

\textbf{cross elasticity of demand} 需求交叉弹性 (XED) measures how much demand for one good responds when the price of \emph{another} good changes.

\[
\text{XED} = \frac{\%\ \text{change in quantity demanded of A}}{\%\ \text{change in price of B}}
\]

\begin{itemize}
\item for \textbf{substitutes}, XED is positive (the price of B rises, so people buy more A).

\item for \textbf{complements}, XED is negative (the price of B rises, so people buy less A and less B).

\item for goods that are not related, XED is about zero.

\end{itemize}

This helps a firm watch its rivals: if XED with a rival's good is high and positive, a price cut by the rival will hurt its sales.

\section*{Price elasticity of supply}

\textbf{price elasticity of supply} 供给价格弹性 (PES) measures how much the quantity supplied responds when the price changes.

\[
\text{PES} = \frac{\%\ \text{change in quantity supplied}}{\%\ \text{change in price}}
\]

PES is positive. Supply is elastic if it responds a lot, inelastic if it responds little. What makes supply more elastic:

\begin{itemize}
\item \textbf{spare capacity} 闲置产能 — unused machines and workers, so firms can quickly make more.

\item \textbf{stocks} 库存 — goods stored in a warehouse that can be sold fast.

\item how easily factors of production can move into the industry.

\item time. Economists split time into three periods:

\begin{itemize}
\item the \textbf{momentary} 瞬时 period — supply is fixed; PES is zero.

\item the \textbf{short run} 短期 — at least one factor is fixed, so supply can rise only a little.

\item the \textbf{long run} 长期 — all factors can change, so supply is most elastic.

\end{itemize}

\end{itemize}

\section*{The interaction of demand and supply}

\subsection*{Equilibrium}

The market is in \textbf{equilibrium} 均衡 where the demand curve crosses the supply curve. Here demand equals supply. This gives the \textbf{equilibrium price} 均衡价格 and the \textbf{equilibrium quantity} 均衡数量. There is no pressure for the price to change.

\subsection*{Disequilibrium}

When the price is not at equilibrium, the market is in \textbf{disequilibrium} 非均衡:

\begin{itemize}
\item if the price is \textbf{too low}, there is \textbf{excess demand} 超额需求 (a \textbf{shortage} 短缺). Buyers compete, so the price is pushed up.

\item if the price is \textbf{too high}, there is \textbf{excess supply} 超额供给. Sellers cannot sell everything, so the price is pushed down.

\end{itemize}

In both cases the price moves back to equilibrium. This shows the price mechanism doing its jobs of signalling, rationing and giving an incentive.

\subsection*{How shifts change the equilibrium}

\begin{center}
\adjustbox{max width=\linewidth}{%
\begin{tabular}{lll}
\toprule
\textbf{Change} & \textbf{Effect on price} & \textbf{Effect on quantity} \\
\midrule
demand rises (shifts right) & rises & rises \\
demand falls (shifts left) & falls & falls \\
supply rises (shifts right) & falls & rises \\
supply falls (shifts left) & rises & falls \\
\bottomrule
\end{tabular}}
\end{center}

\subsection*{Demand and supply in different markets}

The same analysis works in many markets:

\begin{itemize}
\item in \textbf{product markets} 产品市场, demand and supply set the price of goods and services.

\item in \textbf{factor markets} 要素市场, they set the price of factors — for example, the demand for and supply of labour set the \textbf{wage} 工资.

\item in the \textbf{foreign exchange market} 外汇市场, the demand for and supply of a currency set the \textbf{exchange rate} 汇率 (the price of one currency in terms of another).

\end{itemize}

\section*{Consumer and producer surplus}

\textbf{consumer surplus} 消费者剩余 is the gain to buyers. It is the difference between the most a consumer was willing to pay and the price actually paid. On a diagram it is the area below the demand curve and above the price line.

\textbf{producer surplus} 生产者剩余 is the gain to sellers. It is the difference between the price a producer receives and the lowest price the producer would have accepted. On a diagram it is the area above the supply curve and below the price line.

How they change:

\begin{itemize}
\item if demand rises, the price and quantity rise, so producer surplus rises and consumer surplus usually rises too.

\item if supply rises, the price falls and quantity rises, so consumer surplus rises.

\item a higher price (with the curves fixed) raises producer surplus but lowers consumer surplus.

\end{itemize}

Together, consumer surplus plus producer surplus is the total \textbf{welfare} 福利 (gain to society) from the market. It is largest at the free-market equilibrium.


\end{document}
